% ============================================================
%  Příprava na maturitu – Mechatronika
%  Kompiluj: pdflatex maturita_mechatronika.tex
%  Kódování: UTF-8
% ============================================================

\documentclass[a4paper,10pt]{article}

% --- Balíčky ---
\usepackage[T1]{fontenc}
\usepackage[utf8]{inputenc}
\usepackage[czech]{babel}
\usepackage[left=2.4cm, right=1.8cm, top=1.8cm, bottom=1.8cm]{geometry}
\usepackage{lmodern}
\usepackage{microtype}
\usepackage{parskip}
\usepackage{titlesec}
\usepackage{enumitem}
\usepackage{xcolor}
\usepackage{hyperref}
\usepackage{tabularx}
\usepackage{booktabs}
\usepackage{array}
\usepackage{tikz}
\usepackage{eso-pic}
\usepackage{mdframed}
\usepackage{amsmath}
\usepackage{amssymb}
\usepackage{pgfplots}
\pgfplotsset{compat=1.18}
\usetikzlibrary{arrows.meta, positioning, shapes.geometric, calc}

% --- Barvy ---
\definecolor{accent}{HTML}{03254E}
\definecolor{stripeblue}{HTML}{568EA3}
\definecolor{medgray}{HTML}{888888}
\definecolor{lightblue}{HTML}{EAF2F8}
\definecolor{lightyellow}{HTML}{FDFBE4}
\definecolor{lightgreen}{HTML}{EAF7EA}

% --- Modrý pruh vlevo ---
\AddToShipoutPictureBG{%
  \begin{tikzpicture}[remember picture, overlay]
    \fill[stripeblue]
      (current page.north west)
      rectangle
      ([xshift=10mm] current page.south west);
    \fill[accent!60]
      ([xshift=10mm] current page.north west)
      rectangle
      ([xshift=12mm] current page.south west);
  \end{tikzpicture}%
}

\hypersetup{colorlinks=true, urlcolor=accent, linkcolor=accent}

\titleformat{\section}
  {\large\bfseries\color{accent}}
  {\thesection.}
  {0.5em}
  {}
  [\color{accent}\titlerule]

\titlespacing{\section}{0pt}{16pt}{6pt}

\newmdenv[
  backgroundcolor=lightblue,
  linecolor=stripeblue,
  linewidth=1pt,
  roundcorner=4pt,
  innerleftmargin=8pt,
  innerrightmargin=8pt,
  innertopmargin=6pt,
  innerbottommargin=6pt,
  skipabove=4pt,
  skipbelow=4pt
]{pojmybox}

\newmdenv[
  backgroundcolor=lightyellow,
  linecolor=accent!40,
  linewidth=1pt,
  roundcorner=4pt,
  innerleftmargin=8pt,
  innerrightmargin=8pt,
  innertopmargin=6pt,
  innerbottommargin=6pt,
  skipabove=4pt,
  skipbelow=8pt
]{vysvetlenibox}

\newmdenv[
  backgroundcolor=lightgreen,
  linecolor=accent!30,
  linewidth=1pt,
  roundcorner=4pt,
  innerleftmargin=8pt,
  innerrightmargin=8pt,
  innertopmargin=6pt,
  innerbottommargin=6pt,
  skipabove=4pt,
  skipbelow=8pt
]{vzorcebox}

\pagestyle{plain}

% ============================================================
\begin{document}

% --- Záhlaví ---
\begin{minipage}[t]{0.75\textwidth}
    {\Huge\bfseries\color{accent} Příprava na maturitu}\\[4pt]
    {\large\color{stripeblue} Mechatronika}\\[2pt]
    \textcolor{medgray}{\small Suchánek Lukáš \quad SPŠ Prosek \quad 2026}
\end{minipage}
\hfill
\begin{minipage}[t]{0.22\textwidth}
    \raggedleft
    \begin{tikzpicture}
        \draw[color=stripeblue, rounded corners=4pt, line width=1pt]
            (0,0) rectangle (3,1.2)
            node[midway, align=center, color=accent, font=\small\bfseries]
            {Otázek: 22};
    \end{tikzpicture}
\end{minipage}

\vspace{6pt}
\noindent\color{accent}\rule{\linewidth}{1.5pt}
\vspace{2pt}

% ============================================================
\section{Programování v jazyce C}

\begin{pojmybox}
\textbf{Klíčové pojmy:}
datové typy, proměnná, pole, ukazatel, funkce, rekurze,
struktura, preprocesor, dynamická paměť, kompilace, makefile,
zásobník (stack), halda (heap)
\end{pojmybox}

\begin{vysvetlenibox}
\textbf{Datové typy a paměť}

\begin{center}
\begin{tabular}{@{} l l l l @{}}
    \toprule
    \textbf{Typ} & \textbf{Velikost} & \textbf{Rozsah} & \textbf{Použití} \\
    \midrule
    \texttt{uint8\_t}  & 1\,B & 0 -- 255              & registry MCU \\
    \texttt{int16\_t}  & 2\,B & $-32768$ -- $32767$   & rychlé výpočty \\
    \texttt{int32\_t}  & 4\,B & $\pm 2 \times 10^9$   & čítače, čas \\
    \texttt{float}     & 4\,B & $\pm 3{,}4\times10^{38}$ & měření, regulace \\
    \texttt{double}    & 8\,B & $\pm 1{,}8\times10^{308}$ & přesné výpočty \\
    \bottomrule
\end{tabular}
\end{center}

\textbf{Ukazatele -- princip}

\begin{vzorcebox}
\begin{verbatim}
int x = 42;
int *p = &x;   // p obsahuje adresu proměnné x
*p = 100;      // dereferencing -- změní hodnotu x na 100
printf("%d", x); // vypíše 100
\end{verbatim}
\end{vzorcebox}

\textbf{Dynamická paměť (heap)}

\begin{vzorcebox}
\begin{verbatim}
int *pole = malloc(10 * sizeof(int));  // alokace 10 intů
if (pole == NULL) { /* chyba! */ }
pole[0] = 5;
free(pole);   // VŽDY uvolnit, jinak memory leak!
\end{verbatim}
\end{vzorcebox}

\textbf{Struktura paměti programu}

\begin{center}
\begin{tikzpicture}[scale=0.7,
  blk/.style={rectangle, draw=accent, fill=lightblue,
              minimum width=3cm, minimum height=0.6cm,
              text centered, font=\small}]
  \node[blk, fill=lightgreen]  at (0, 3)   {Stack (zásobník)};
  \node[blk, fill=lightyellow] at (0, 2.2) {Heap (halda)};
  \node[blk]                   at (0, 1.4) {BSS (neinic. data)};
  \node[blk]                   at (0, 0.6) {Data (init. data)};
  \node[blk, fill=lightblue]   at (0,-0.2) {Text (kód programu)};
  \draw[>-, accent] (-1.8,3) -- (-1.8,2.4)
      node[midway,left,font=\scriptsize]{roste dolů};
  \draw[>-, stripeblue] (-1.8,2) -- (-1.8,2.6)
      node[midway,left,font=\scriptsize]{roste nahoru};
\end{tikzpicture}
\end{center}

\textbf{Preprocesor -- direktiva \texttt{\#define}}\\
Makra se nahrazují \textit{před} kompilací -- nulová režie za běhu:
\begin{vzorcebox}
\begin{verbatim}
#define PI        3.14159f
#define MAX(a,b)  ((a) > (b) ? (a) : (b))
#define LED_PIN   13
\end{verbatim}
\end{vzorcebox}
\end{vysvetlenibox}

% ============================================================

% ============================================================
\end{document}
